\chapter{Plasmapheresis}

\section*{What Is This Test or Procedure?}
Plasmapheresis, or therapeutic plasma exchange, is a procedure that removes plasma containing pathogenic antibodies, immune complexes, and toxins from the blood and replaces it with donor plasma or albumin. The cellular components of blood are returned to the patient.

\section*{Alternative Names}
Therapeutic plasma exchange; TPE; plasma exchange; apheresis

\section*{Principle}
Blood is withdrawn through a venous catheter and passed through a centrifuge or membrane filter that separates plasma from the cellular elements. The plasma is discarded and replaced with colloid solution, and the cellular components are returned to the circulation, selectively removing disease-causing substances from the blood.

\section*{History}
Therapeutic apheresis was pioneered in the 1960s and 1970s with the development of continuous-flow centrifugal cell separators, and it became a standard treatment for antibody-mediated and immune-mediated disorders.

\section*{Why This Test or Procedure Is Ordered}
Ordered to treat conditions caused or worsened by pathogenic antibodies, immune complexes, or paraproteins, including Guillain-Barr\'e syndrome, myasthenia gravis crisis, thrombotic thrombocytopenic purpura (TTP), antibody-mediated rejection of transplanted organs, antiglomerular basement membrane disease, and severe hyperviscosity syndromes.

\section*{Is This a Primary or Secondary Test or Procedure}
Primary treatment for selected immune-mediated and antibody-mediated disorders; it is a therapeutic rather than a diagnostic procedure.

\section*{How This Test or Procedure Is Performed}
Vascular access is obtained through a large-bore central or peripheral catheter. Blood is drawn into an apheresis machine that separates plasma from cellular components by centrifugation or filtration. The plasma is replaced with albumin or fresh frozen plasma, and the cellular components are reinfused; treatment typically lasts one to two hours.

\section*{What Preparation Is Needed}
Informed consent, review of the coagulation profile and platelet count, confirmation of adequate vascular access, and baseline laboratory studies are required. Anticoagulant (citrate) is used in the circuit, and the patient should be aware of the required number of sessions.

\section*{Contraindications and Precautions}
Hemodynamic instability, uncontrolled sepsis, inability to obtain vascular access, and sensitivity to replacement fluids are contraindications. Precautions include hypocalcemia from citrate toxicity, hypotension, and use of fresh frozen plasma in patients with coagulation abnormalities or risk of bleeding.

\section*{Risks}
Risks include citrate-induced hypocalcemia, hypotension, allergic reactions to plasma, bleeding, catheter-related infection or thrombosis, and rarely fluid overload.

\section*{What Happens After the Test or Procedure}
The patient is monitored for hypotension, hypocalcemia, and bleeding before discharge or return to the ward. Serial clinical and laboratory monitoring, such as platelet counts in TTP, determines the response and the need for additional sessions.

\section*{Understanding Your Results}
Clinical response is judged by improvement in symptoms and organ function, with laboratory markers such as ADAMTS13 activity, platelet count, antibody titers, or paraprotein levels monitored. A partial or absent response may prompt additional sessions or alternative therapy.

\section*{Normal Results}
Return of disease-specific laboratory markers toward normal (for example, rising platelet count in TTP, falling antibody titers) indicates an adequate therapeutic response.

\section*{Follow-up Tests and Procedures}
Continue immunosuppressive or disease-specific therapy that addresses the underlying disorder, since plasmapheresis alone does not stop antibody production. Serial laboratory monitoring and repeat treatment courses may be required for relapsing disease.

\section*{References}
\begin{enumerate}
  \item MedlinePlus. Plasmapheresis. U.S. National Library of Medicine; 2025.
  \item Current Medical Diagnosis and Treatment. McGraw-Hill; 2025.
\end{enumerate}

\clearpage
